\documentclass[12pt]{article}
\usepackage[utf8]{inputenc}
\usepackage[T1]{fontenc}
\usepackage{hyperref}
\title{Phola Language Resources for K-12 Curriculum Development}
\author{Compiled by Claude Sonnet 5.5 (claude-sonnet-5-5) for Andr\'as Kornai}
\date{}
\begin{document}
\maketitle

\section{Introduction}
Phola (ISO 639-3 ypg, Glottocode phol1237; also called Upriver Phula) is a dialect cluster of the Loloish languages, spoken by the Phula people of Yunnan, China. Wikipedia gives about 13,000 speakers for Phola proper, in Honghe, Gejiu, Shiping, Jianshui and Yuanjiang counties. The brief gives about 25,000, and the two figures differ. The best-studied variety is that of Luodie, Yuanjiang County (Pelkey 2011). Phola belongs to the Riverine group of Southeastern Loloish. It should not be confused with Pholo (yip), which is unrelated. The code ypg also covers Phala (ypa) and Alo Phola (ypo), which share the Phola name.

\section{Primary Educational Resources}
No primers, textbooks or school programmes for Phola turned up in the searches. The language is not known to be taught in schools.

\section{Linguistic and Reference Materials}
\begin{itemize}
\item Phola language (Wikipedia): \url{https://en.wikipedia.org/wiki/Phola_language}
\item Pelkey, ``Puzzling over Phula'' (academic survey): \url{https://www.academia.edu/2435607/Puzzling_over_Phula}
\item Southeastern Loloish languages (Wikipedia, classification context): \url{https://en.wikipedia.org/wiki/Southeastern_Loloish_languages}
\item Loloish languages (Wikipedia, broader context): \url{https://en.wikipedia.org/wiki/Loloish_languages}
\item Pholo language (Wikipedia; a different language, listed to avoid confusion): \url{https://en.wikipedia.org/wiki/Pholo_language}
\end{itemize}

\section{Religious and Community Resources}
\begin{itemize}
\item Joshua Project profile of Phola in China (an ethnic-group profile; I did not check its contents beyond the title): \url{https://joshuaproject.net/index.php/people_groups/16151}
\end{itemize}

\section{Digital and Audio Resources}
None were found for Phola in the searches.

\section{Curriculum Development Context}
Searches found no orthography, dictionary, grammar or classroom materials for Phola. Pelkey's work on the Luodie variety is the best-documented base. A curriculum would need new community-led documentation and a decision about which variety to use as the standard.
\end{document}